% =====================================================================
\section{Experimentos}\label{sec:exp}
% =====================================================================

Evaluamos Walk Attention sobre una familia controlada de grafos con
over-squashing severo y \emph{ajustable}, para aislar la contribución de cada
componente arquitectónico. Publicamos junto al artículo todo el código, las
configuraciones y los registros de los experimentos.

% ---------------------------------------------------------------------
\subsection{Un benchmark de over-squashing ajustable}
% ---------------------------------------------------------------------

\paragraph{Grafos de cadena con cuello de botella.}
Construimos una familia parametrizada por una \emph{profundidad} $d$, un número de fuentes $K$
y un ancho $M$. El grafo tiene $K$ vértices fuente, luego $d$ \emph{capas de cuello de botella}
consecutivas de $M$ vértices intercambiables cada una, y un
vértice objetivo. Cada fuente se conecta con cada vértice de la primera capa, cada
vértice de la capa $r$ con cada vértice de la capa $r+1$, y cada vértice de
la última capa con el objetivo. Como los vértices del cuello de botella son
intercambiables, hay $M^{d}$ caminos de longitud $(d{+}1)$ desde una fuente al objetivo,
de modo que la multiplicidad de caminos hacia el objetivo crece como
$n_{d+1}(\cdot,\text{target})\sim K\,M^{d}$
(Proposición~\ref{prop:count}); una GNN de paso de mensajes debe comprimirlos todos
en el vector de ancho fijo del objetivo. El soporte de alcanzabilidad por caminos, en cambio,
permanece disperso: en el rango más profundo es solo el de los $K$ pares
fuente--objetivo ($5/196\approx2.6\%$ de todos los pares de nodos a profundidad
$2$ y $5/324\approx1.5\%$ a profundidad $3$ en nuestra configuración).

\paragraph{Tarea.}
Cada fuente lleva una \emph{clave} one-hot y una \emph{carga útil} escalar; el objetivo
lleva una clave de consulta. El modelo debe producir en el objetivo la clave de la
fuente que coincide con la consulta---una recuperación cuya respuesta se sitúa a $d{+}1$
saltos, detrás del cuello de botella. La precisión se mide en el objetivo;
el azar es $1/K$.

\paragraph{Modelos.} Comparamos:
\begin{itemize}
  \item \textbf{GAT} \cite{velickovic2017graph}: atención de grafos de un salto, $d{+}1$
  capas (pesos aprendidos, soporte de un salto);
  \item \textbf{Operador de caminos} (\texttt{walkraw}): el operador multi-salto de pesos fijos
  $W^{(k)}$ de la Definición~\ref{def:walkop} (soporte multi-salto, pesos fijos
  de multiplicidad);
  \item \textbf{Walk Attention} (nuestro): Algoritmo~\ref{alg:wa} (soporte multi-salto,
  pesos aprendidos).
\end{itemize}
Todos los modelos usan ancho oculto $16$, cuatro cabezas cuando corresponde, y $d{+}1$
capas/rangos. Entrenamos con AdamW, warmup lineal de diez épocas
seguido de decaimiento coseno, y recorte de gradiente; Walk Attention añade
normalización de capa tras cada capa. Reportamos media y desviación estándar de
la precisión en el objetivo sobre una partición de validación reservada (la tarea
sintética no tiene partición de prueba separada) y cinco semillas, en
profundidades $d\in\{2,3\}$ con $K=5$, $M=4$.

% ---------------------------------------------------------------------
\subsection{Resultado principal}
% ---------------------------------------------------------------------

La Tabla~\ref{tab:main} reporta los resultados, que caen exactamente donde las
Proposiciones~\ref{prop:collapse}--\ref{prop:wa} los sitúan
(Observación~\ref{rem:mechanisms}). La GAT de un salto se degrada
abruptamente con la profundidad, acercándose al azar: un operador local no puede
propagar la señal a través del cuello de botella bajo over-squashing. GCN y GIN,
cuya primera agregación es lineal, quedan clavados en el azar por la
Proposición~\ref{prop:blind}---y allí se quedan empíricamente a toda
profundidad. El operador de caminos de pesos fijos alcanza
el objetivo---su soporte multi-salto cubre el cuello de botella en una sola capa---pero, al ponderar todas
las fuentes alcanzables solo por multiplicidad, no puede seleccionar la que coincide y
se estanca en $0.50$--$0.62$. \textbf{Walk Attention resuelve la tarea exactamente},
con precisión $1.000\pm0.000$ en ambas profundidades y en las cinco semillas: el mismo
soporte multi-salto, ahora con pesos consulta--clave aprendidos, permite que el objetivo atienda
selectivamente a la fuente relevante.

\begin{table}[t]
\centering
\caption{Precisión en el objetivo (media $\pm$ desv.\ estándar sobre 5 semillas) en la recuperación
en cadena con cuello de botella, $K{=}5$, $M{=}4$. El azar es $1/K=0.20$. La profundidad $d$ es el número de
capas de cuello de botella; la respuesta se sitúa a $d{+}1$ saltos de distancia.
GCN y GIN quedan en el azar, como exige la Proposición~\ref{prop:blind} (el
residuo de $+0.01$--$0.02$ sobre $1/K$ es el sesgo de selección de mejor época
sobre una partición de validación finita).}
\label{tab:main}
\begin{tabular}{lcc}
\toprule
\textbf{Model} & \textbf{depth $d=2$} & \textbf{depth $d=3$} \\
\midrule
GCN (one-hop, linear aggregation)  & $0.21\pm0.01$ & $0.22\pm0.01$ \\
GIN (one-hop, linear aggregation)  & $0.21\pm0.01$ & $0.21\pm0.01$ \\
GAT (one-hop attention)            & $0.45\pm0.22$ & $0.29\pm0.17$ \\
Walk operator (fixed weights)      & $0.62\pm0.12$ & $0.50\pm0.12$ \\
\textbf{Walk Attention (ours)}     & $\mathbf{1.000\pm0.000}$ & $\mathbf{1.000\pm0.000}$ \\
\bottomrule
\end{tabular}
\end{table}

\paragraph{Sobre la estabilidad.}
Sin normalización de capa ni warmup, Walk Attention tiene un techo alto pero
alta varianza: su precisión osciló entre $0.44$ y $1.00$ según la semilla. Los tres
estabilizadores estándar (normalización de capa, warmup, recorte de gradiente)
eliminan por completo esta varianza, dando el $\pm0.000$ de la
Tabla~\ref{tab:main}. La inestabilidad fue, pues, un artefacto de optimización, no
una propiedad del operador.

% ---------------------------------------------------------------------
\subsection{RingTransfer}
% ---------------------------------------------------------------------

La familia de cuellos de botella es construcción propia. Para confirmar el efecto
en un benchmark \emph{independiente} usamos \textsc{RingTransfer}
\cite{digiovanni2023oversquashing}: un ciclo de $n$ nodos donde una fuente debe
transmitir su etiqueta one-hot al objetivo diametralmente opuesto. La estructura de
caminos paralelos es explícita: la fuente alcanza el objetivo por los \emph{dos arcos}
del anillo, dos caminos de igual longitud $n/2$. El over-squashing es severo y
ajustable: cada nodo tiene dos vecinos, así que el número de recorridos de
longitud $n/2$ que entran al objetivo---los mensajes que una red de un salto debe
comprimir tras $n/2$ capas---crece como $2^{\,n/2}$ con el tamaño del anillo,
mientras solo los dos recorridos por los arcos portan la señal. La tarea requiere
$n/2$ saltos, y por tanto $n/2$ capas en un modelo de un salto. Usamos cinco
clases (azar $0.2$), ancho oculto $32$, la misma receta de optimización (AdamW,
warmup lineal, decaimiento coseno, recorte de gradiente), y $n/2$ capas/rangos
para todos los modelos.

La Tabla~\ref{tab:ring} reporta los resultados. En anillos pequeños
($n\le 10$) todos los modelos tienen éxito, pues la compresión aún no es severa.
Desde $n=14$ (distancia $7$) las líneas base de un salto se degradan y se vuelven
muy dependientes de la semilla---en $n=18$ GAT cae a $0.29\pm0.19$ y GCN a
$0.53\pm0.33$---mientras que \textbf{Walk Attention se mantiene en
$1.000\pm0.000$}: su soporte multi-salto agrega la fuente directamente por
cualquiera de los dos arcos, y los pesos aprendidos identifican la etiquetada.
La brecha se amplía con el tamaño del anillo, esto es, con la severidad del
over-squashing, como se predijo.

\begin{table}[t]
\centering
\caption{Precisión de \textsc{RingTransfer} frente al tamaño del anillo $n$
(distancia $n/2$), media $\pm$ desv.\ estándar sobre 5 semillas; azar $=0.20$.
GCN/GAT de un salto colapsan a medida que el anillo---y por tanto el
over-squashing---crece, mientras que Walk Attention se mantiene exacto.}
\label{tab:ring}
\begin{tabular}{lcccc}
\toprule
\textbf{Model} & $n{=}6$ & $n{=}10$ & $n{=}14$ & $n{=}18$ \\
\midrule
GCN                    & $1.00\pm0.00$ & $1.00\pm0.00$ & $0.81\pm0.13$ & $0.53\pm0.33$ \\
GAT                    & $1.00\pm0.00$ & $1.00\pm0.00$ & $0.69\pm0.29$ & $0.29\pm0.19$ \\
\textbf{Walk Attention (ours)} & $\mathbf{1.00\pm0.00}$ & $\mathbf{1.00\pm0.00}$ & $\mathbf{1.00\pm0.00}$ & $\mathbf{1.00\pm0.00}$ \\
\bottomrule
\end{tabular}
\end{table}

% ---------------------------------------------------------------------
\subsection{Colapsar caminos perjudica}
% ---------------------------------------------------------------------

El cociente $\kQ/I$ de la Sección~\ref{sec:algebra} sugiere un
remedio alternativo: reemplazar las multiplicidades $n_k(i,j)$ por las dimensiones
menores del cociente $\dim(e_j(\kQ/I)_k e_i)$, deduplicando caminos equivalentes
antes de la agregación. Implementamos esto (el operador de pesos fijos construido a partir de las
dimensiones del cociente) y lo comparamos con el operador sin cociente sobre la
misma tarea, en una ablación controlada que cambia \emph{únicamente} el
operador (cociente frente a crudo) y mantiene todo lo demás---datos,
arquitectura y la receta de optimización de la Tabla~\ref{tab:main}---fijo. Como
muestra la Tabla~\ref{tab:quotient}, \textbf{el operador del cociente cae
claramente por debajo del crudo en ambas profundidades}. Colapsar caminos
paralelos descarta la multiplicidad de fuentes que la recuperación necesita: en
este régimen los caminos redundantes son \emph{señal}, no ruido, y eliminarlos
elimina información.

\begin{table}[t]
\centering
\caption{Deduplicar caminos vía el cociente $\kQ/I$ \emph{perjudica}: precisión
en el objetivo (media $\pm$ desv.\ estándar sobre 5 semillas), recuperación con
cuello de botella, bajo el mismo protocolo de la Tabla~\ref{tab:main} (la fila
del operador de caminos coincide con ella). El operador del cociente queda
claramente por debajo del operador sin cociente en ambas profundidades.}
\label{tab:quotient}
\begin{tabular}{lcc}
\toprule
\textbf{Operator} & \textbf{depth $d=2$} & \textbf{depth $d=3$} \\
\midrule
Quotient $\kQ/I$ (collapsed walks) & $0.39\pm0.04$ & $0.30\pm0.09$ \\
Walk operator (raw walks)          & $\mathbf{0.62\pm0.12}$ & $\mathbf{0.50\pm0.12}$ \\
\bottomrule
\end{tabular}
\end{table}

Este resultado negativo indica que la contribución del álgebra de caminos
aquí \emph{no} es la compresión. Su papel apropiado es proporcionar
el \emph{soporte} de la atención multi-salto---disperso y exacto---sobre
el cual la red aprende los pesos, que es precisamente Walk
Attention. Si existe un régimen en el que la compresión del cociente sea
beneficiosa (p.\ ej.\ cuando los caminos paralelos son ruido redundante, o
cuando clases de caminos distintas se enrutan a cabezas distintas) queda
abierto.

% ---------------------------------------------------------------------
\subsection{LRGB Peptides-func}
% ---------------------------------------------------------------------

El benchmark sintético es deliberadamente un peor caso: un cuello de botella extremo donde
el alcance multi-salto es decisivo. Para evaluar Walk Attention con datos reales
usamos también \textbf{Peptides-func} del Long Range Graph Benchmark
\cite{dwivedi2022lrgb} ($10{,}873/2{,}331/2{,}331$ grafos moleculares de entrenamiento/validación/prueba,
mediana $138$ nodos; clasificación multietiqueta de grafos, métrica Average Precision).
Usamos el mismo embedding de átomos, cuatro capas, ancho oculto $80$, pooling de media
y AdamW para todos los modelos; Walk Attention usa los rangos $k=1,\dots,4$ con
el soporte de alcanzabilidad por caminos de cada molécula. Aquí el soporte sigue
disperso ($1$--$8\%$ de $n^2$ en el rango $3$), de modo que el método sigue barato.

En AP de prueba obtenemos GCN $0.492\pm0.011$, GAT $0.526\pm0.004$ y Walk
Attention $0.526\pm0.007$ (media $\pm$ desv.\ estándar sobre cuatro semillas).
\textbf{Walk Attention es comparable a las líneas base de un salto pero no las
supera aquí}: supera claramente a GCN y es estadísticamente indistinguible de
GAT. Dos salvedades enmarcan esta comparación. Primera, todos los modelos comparten un
presupuesto deliberadamente pequeño (cuatro capas, ancho $80$, sin
codificaciones posicionales, schedule corto), así que nuestros números absolutos
están por debajo de los publicados para este benchmark---el artículo de LRGB
reporta GCN en $0.593$ bajo su protocolo afinado de $500$k parámetros
\cite{dwivedi2022lrgb}, y re-evaluaciones posteriores son aún mayores
\cite{tonshoff2023where}; la comparación es por tanto \emph{interna}, aislando
el operador de agregación bajo un protocolo idéntico, no una afirmación frente
al estado del arte. Segunda, el resultado no se debe a la ausencia de caminos
paralelos: los
grafos de péptidos tienen anillos y ramificaciones, y una molécula típica
tiene cientos de pares de vértices unidos por $\ge 2$ caminos de longitud $\le 4$
(que el ideal del cociente colapsaría). Más bien, el over-squashing es
\emph{leve}: estos grafos son pequeños (mediana $138$ nodos), así que pocas
capas de atención de un salto ya alcanzan lo suficiente, y la
multiplicidad hacia cualquier nodo no crece lo bastante para saturar el
embedding. El soporte multi-salto, decisivo bajo over-squashing severo, aporta
poco cuando la compresión es leve. La ventaja de Walk Attention es por tanto
\emph{dependiente del régimen}---grande bajo compresión severa
(Tablas~\ref{tab:main} y~\ref{tab:ring}) y comparable a la atención de un salto cuando es leve---y,
con datos reales, el operador sigue entrenando de forma estable a costo moderado
(alrededor de $1.6\times$ el tiempo de pared de la GAT de un salto, promediado
sobre semillas, muy por debajo de la atención densa de todos los pares).

% ---------------------------------------------------------------------
\subsection{Costo}
% ---------------------------------------------------------------------

Walk Attention opera sobre el soporte de alcanzabilidad por caminos, de tamaño el
número de entradas no nulas de $\Adj^{\,k}$: una pequeña fracción de $n^2$
($1.5$--$2.6\%$ en el rango más profundo, según la profundidad), por lo que la
capa es mucho más barata que
un Transformer denso sin perder largo alcance. Los soportes se
precalculan y cachean por topología, fuera del bucle de entrenamiento.
